%; whizzy paragraph -pdf xpdf -latex ./whizzypdfptex.sh
%; whizzy-paragraph "^\\\\begin{frame}\\|\\\\emtext"
% latex beamer presentation.
% platex, latex-beamer でコンパイルすることを想定。 

%     Tokyo Debian Meeting resources
%     Copyright (C) 2012 Junichi Uekawa

%     This program is free software; you can redistribute it and/or modify
%     it under the terms of the GNU General Public License as published by
%     the Free Software Foundation; either version 2 of the License, or
%     (at your option) any later version.

%     This program is distributed in the hope that it will be useful,
%     but WITHOUT ANY WARRANTY; without even the implied warreanty of
%     MERCHANTABILITY or FITNESS FOR A PARTICULAR PURPOSE.  See the
%     GNU General Public License for more details.
%     You should have received a copy of the GNU General Public License
%     along with this program; if not, write to the Free Software
%     Foundation, Inc., 51 Franklin St, Fifth Floor, Boston, MA  02110-1301 USA

\documentclass[cjk,dvipdfmx,12pt]{beamer}
\usetheme{Tokyo}
\usepackage{monthlypresentation}

%  preview (shell-command (concat "evince " (replace-regexp-in-string "tex$" "pdf"(buffer-file-name)) "&")) 
%  presentation (shell-command (concat "xpdf -fullscreen " (replace-regexp-in-string "tex$" "pdf"(buffer-file-name)) "&"))
%  presentation (shell-command (concat "evince " (replace-regexp-in-string "tex$" "pdf"(buffer-file-name)) "&"))

%http://www.naney.org/diki/dk/hyperref.html
%日本語EUC系環境の時
\AtBeginDvi{\special{pdf:tounicode EUC-UCS2}}
%シフトJIS系環境の時
%\AtBeginDvi{\special{pdf:tounicode 90ms-RKSJ-UCS2}}

\newenvironment{commandlinesmall}%
{\VerbatimEnvironment
  \begin{Sbox}\begin{minipage}{1.0\hsize}\begin{fontsize}{8}{8} \begin{BVerbatim}}%
{\end{BVerbatim}\end{fontsize}\end{minipage}\end{Sbox}
  \setlength{\fboxsep}{8pt}
% start on a new paragraph

\vspace{6pt}% skip before
\fcolorbox{dancerdarkblue}{dancerlightblue}{\TheSbox}

\vspace{6pt}% skip after
}
%end of commandlinesmall

\title{東京エリアDebian勉強会}
\subtitle{第147回 2017年1月度}
\author{杉本 典充  dictoss@live.jp}
\date{2017年1月21日}
\logo{\includegraphics[width=8cm]{image200607/openlogo-light.eps}}

\begin{document}

\begin{frame}
\titlepage{}
\end{frame}

\begin{frame}{Agenda}
 \begin{minipage}[t]{0.45\hsize}
  \begin{itemize}
  \item 最近あったDebian関連のイベント報告
	\begin{itemize}
	\item 第146回東京エリアDebian勉強会
    \item Mini Debian Conference Japan 2016
	\end{itemize}
  \item 事前課題発表
  \end{itemize}
 \end{minipage}
 \begin{minipage}[t]{0.45\hsize}
  \begin{itemize}
   \item Debian Trivia Quiz
  \end{itemize}
 \end{minipage}
\end{frame}

\section{イベント報告}
\emtext{イベント報告}

\begin{frame}{第146回東京エリアDebian勉強会}

\begin{itemize}
\item 場所は朝日ネットさんをお借りして開催しました
\item 参加者は5名
\item セミナ内容は、以下内容を発表しました
  \begin{itemize}
  \item 吉野さん「dh\_strip\_nondeterminismについて」
  \end{itemize}
\item 残りの時間でhack timeを行い、成果発表をしました
\item \url{https://tokyodebianmeeting.titanpad.com/ep/pad/export/146/latest?format=txt}
\end{itemize} 
\end{frame}

\begin{frame}{Mini Debian Conference Japan 2016}

\begin{itemize}
\item 12/10(土)にサイボウズ様で開催されました
\item 「LibreOffice Kaigi 2016.12」と併催
\item 国内のdebian開発者やdebianユーザが集まりました
\item 台湾などからDebian Developerの方々も参加しました
\item タイムテーブルは以下webサイトを参照ください
\item \url{http://miniconf.debian.or.jp/}
\end{itemize} 
\end{frame}

\section{事前課題}
\emtext{事前課題}
{\footnotesize
 \begin{prework}{ dictoss }
  \begin{enumerate}
  \item stretch$B$N%P%0$D$V$7(B
  \end{enumerate}
\end{prework}

\begin{prework}{ yy\_y\_ja\_jp }
  \begin{enumerate}
  \item DDTSS
  \end{enumerate}
\end{prework}

\begin{prework}{ oguraysu }
  \begin{enumerate}
  \item postfix$B$N@_Dj$r$7$F$_$k(B
  \end{enumerate}
\end{prework}

\begin{prework}{ henrich }
  \begin{enumerate}
  \item $B%&%'%V%5%$%H$N:F9=C[(B
  \end{enumerate}
\end{prework}

\begin{prework}{ upsilon }
  \begin{enumerate}
  \item TrimSlice$B$K(BDebian(armhf)$B$r%$%s%9%H!<%k$7$^$9(B
  \end{enumerate}
\end{prework}

\begin{prework}{ kenhys }
  \begin{enumerate}
  \item $B%Q%C%1!<%8$N%a%s%F%J%s%9(B
  \end{enumerate}
\end{prework}

\begin{prework}{ takaswie }
  \begin{enumerate}
  \item $B8=:_3+H/Cf$N%+!<%M%k%b%8%e!<%k(B(snd-firewire-motu)$B$N:n6H$r$7$^$9!#(B
  \end{enumerate}
\end{prework}

}

%\subsection{問題}
%
%%; whizzy-master ../debianmeetingresume201311.tex
% $B0J>e$N@_Dj$r$7$F$$$k$?$a!"$3$N%U%!%$%k$G(B M-x whizzytex $B$9$k$H!"(Bwhizzytex$B$,MxMQ$G$-$^$9!#(B
%

\santaku
{Debian 9 (stretch)$B$,(B 2017-01-07$B$K(Bsoft freeze$B$KF~$C$?$3$H$,%"%J%&%s%9$5$l$^$7$?!#(Bsoft freeze$B$K$J$k$H2?$,$G$-$J$/$J$k$G$7$g$&$+!#(B}
{Debian Wiki$B$N(Bstretch$B%Z!<%8$NJQ99(B}
{$B?75,$N%=!<%9%Q%C%1!<%8$N%"%C%W%m!<%I(B}
{$B4{B8$N%=!<%9%Q%C%1!<%8$N%"%C%W%m!<%I(B}
{B}
{soft freeze$B$KF~$k$H4{B8%Q%C%1!<%8$NIJ<A8~>e$KCmNO$;$h!"$H$$$&$3$H$G?75,$N%=!<%9%Q%C%1!<%8$r%"%C%W%m!<%I$G$-$J$/$J$j$^$9!#(Bfull freeze$B$O(B2017-02-05$B$KM=Dj$5$l$F$$$^$9!#%Q%C%1!<%8$,(Bunstable$B$+$i(Btesting$B$X0\9T$9$k$^$G:GC;(B10$BF|4V$+$+$k$?$a!"Aa$a$K=$@5$7$F%"%C%W%m!<%I$7$^$7$g$&!#(B\url{https://lists.debian.org/debian-devel-announce/2017/01/msg00002.html}}


\section{セミナー1}
\emtext{debhelper 10}

\section{Hack time}
\emtext{Hack time}

\begin{frame}{成果を記入下さい!}

  今回、Hack time時の成果を記録に残してみます。

\url{https://tokyodebianmeeting.titanpad.com/147}

に皆さんアクセス頂き（認証不要です）、各自成果を17:30までに記録するようにお願いします。

\end{frame}
  
\section{今後のイベント}
\emtext{今後のイベント}
\begin{frame}{今後のイベント}
  \begin{itemize}
  \item 2017/2/11(土) 東京エリアDebian勉強会
    \begin{itemize}
    \item BSP(バグスカッシュパーティ)の予定
    \item 時間は調整中ですが、10時位から夕方まで考えています
    \end{itemize}
  \item 2017/3/11(土) OSC 2017 Tokyo/Spring
    \begin{itemize}
    \item セミナーによる出張勉強会、ブースでの展示
    \end{itemize}
  \item 2017/3/18(土) 東京エリアDebian勉強会(実施するかどうか含めて検討中)
  \end{itemize}
\end{frame}

\end{document}

;;; Local Variables: ***
;;; outline-regexp: "\\([ 	]*\\\\\\(documentstyle\\|documentclass\\|emtext\\|section\\|begin{frame}\\)\\*?[ 	]*[[{]\\|[]+\\)" ***
;;; End: ***
